\documentclass[11pt]{article}
% ---------- shared preamble ----------
\usepackage[a4paper,top=26mm,bottom=26mm,left=25mm,right=25mm,headheight=14pt]{geometry}
\usepackage{amsmath,amssymb,amsthm,mathtools}
\usepackage{fontspec}
\setmainfont{Noto Serif CJK SC}[
  BoldFont={Noto Serif CJK SC Bold},
  ItalicFont={Noto Serif CJK SC},ItalicFeatures={FakeSlant=0.15},
  BoldItalicFont={Noto Serif CJK SC Bold},BoldItalicFeatures={FakeSlant=0.15}]
\setsansfont{Noto Sans CJK SC}[BoldFont={Noto Sans CJK SC Bold}]
\newfontfamily\lmr{Latin Modern Roman}
\newcommand{\apos}{{\lmr\char"2019}}
\newcommand{\gerlq}{{\lmr\char"201E}}
\newcommand{\gerrq}{{\lmr\char"201C}}
\XeTeXlinebreaklocale "zh"
\XeTeXlinebreakskip = 0pt plus 1pt
\linespread{1.32}
\setlength{\parindent}{2em}
\setlength{\parskip}{0.35em}
\setcounter{tocdepth}{1}
\emergencystretch=2em

\usepackage{xcolor}
\definecolor{accent}{HTML}{1F4E79}
\definecolor{accentlight}{HTML}{EAF1F8}
\definecolor{warm}{HTML}{8A4B08}
\definecolor{warmlight}{HTML}{FBF3E6}
\definecolor{good}{HTML}{2E6B30}
\definecolor{goodlight}{HTML}{EEF6EE}
\definecolor{muted}{HTML}{5F6B76}

\usepackage{booktabs,array,multirow,enumitem,float}
\setlist{itemsep=0.15em,topsep=0.3em,leftmargin=2.2em}
\usepackage{tikz}
\usetikzlibrary{arrows.meta,calc,positioning,decorations.pathreplacing}
\usepackage{pgfplots}
\pgfplotsset{compat=1.18}
\usepackage[most]{tcolorbox}
\usepackage{titlesec}
\titleformat{\section}{\sffamily\bfseries\Large\color{accent}}{\thesection}{0.8em}{}
\titleformat{\subsection}{\sffamily\bfseries\large\color{accent}}{\thesubsection}{0.7em}{}
\titleformat{\subsubsection}{\sffamily\bfseries\normalsize}{\thesubsubsection}{0.6em}{}
\titlespacing*{\section}{0pt}{1.4em}{0.7em}
\titlespacing*{\subsection}{0pt}{1.1em}{0.5em}

\usepackage{fancyhdr}
\pagestyle{fancy}
\fancyhf{}
\renewcommand{\headrulewidth}{0.4pt}
\fancyfoot[C]{\small\color{muted}\thepage}

\usepackage{hyperref}
\hypersetup{colorlinks=true,linkcolor=accent,citecolor=accent,urlcolor=accent}
\usepackage{caption}
\captionsetup{font=small,labelfont={bf,sf},labelsep=quad}
\renewcommand{\figurename}{图}
\renewcommand{\tablename}{表}
\renewcommand{\contentsname}{目录}
\renewcommand{\refname}{参考文献}
\renewcommand{\proofname}{证明}

% theorem styles (upright body; Chinese should not be slanted)
\newtheoremstyle{zhthm}{0.7em}{0.7em}{\normalfont}{}{\sffamily\bfseries\color{accent}}{}{0.6em}{\thmname{#1}\thmnumber{ #2}\thmnote{（#3）}}
\newtheoremstyle{zhrem}{0.5em}{0.5em}{\normalfont}{}{\sffamily\bfseries\color{muted}}{}{0.6em}{\thmname{#1}\thmnumber{ #2}\thmnote{（#3）}}
\theoremstyle{zhthm}
\newtheorem{theorem}{定理}[section]
\newtheorem{lemma}[theorem]{引理}
\newtheorem{proposition}[theorem]{命题}
\newtheorem{corollary}[theorem]{推论}
\newtheorem{definition}[theorem]{定义}
\newtheorem{assumption}[theorem]{假设}
\theoremstyle{zhrem}
\newtheorem{remark}[theorem]{注}
\newtheorem{example}[theorem]{例}
\makeatletter
\renewenvironment{proof}[1][\proofname]{\par\pushQED{\qed}\normalfont\topsep6\p@\@plus6\p@\relax
  \trivlist\item[\hskip\labelsep\sffamily\bfseries #1.]\ignorespaces}{\popQED\endtrivlist\@endpefalse}
\makeatother

% math macros
\newcommand{\1}{\mathbf{1}}
\newcommand{\Prob}{\mathbb{P}}
\newcommand{\E}{\mathbb{E}}
\newcommand{\R}{\mathbb{R}}
\newcommand{\N}{\mathbb{N}}
\newcommand{\Bin}{\operatorname{Bin}}
\newcommand{\Corr}{\operatorname{Corr}}
\newcommand{\Amin}{A_{-}}
\newcommand{\Aplus}{A_{+}}
\newcommand{\Aind}{A_{\mathrm{ind}}}
\newcommand{\pd}[1]{p_{(#1)}}   % descending order
\newcommand{\pa}[1]{p_{[#1]}}   % ascending order

% boxes
\newtcolorbox{keybox}[1][]{enhanced,breakable,colback=accentlight,colframe=accent,boxrule=0.6pt,arc=2pt,left=8pt,right=8pt,top=6pt,bottom=6pt,fonttitle=\sffamily\bfseries,#1}
\newtcolorbox{warnbox}[1][]{enhanced,breakable,colback=warmlight,colframe=warm,boxrule=0.6pt,arc=2pt,left=8pt,right=8pt,top=6pt,bottom=6pt,fonttitle=\sffamily\bfseries,#1}
\newtcolorbox{goodbox}[1][]{enhanced,breakable,colback=goodlight,colframe=good,boxrule=0.6pt,arc=2pt,left=8pt,right=8pt,top=6pt,bottom=6pt,fonttitle=\sffamily\bfseries,#1}

\newcommand{\doctitle}[3]{%
  \thispagestyle{empty}
  \begin{center}
  {\sffamily\small\color{muted} DSAA2011 Task B \,·\, Project 2 \,·\, Ensemble Voting Beyond Independence}\par\vspace{1.0em}
  {\sffamily\bfseries\LARGE\color{accent} #1}\par\vspace{0.6em}
  {\sffamily\large #2}\par\vspace{0.8em}
  {\small\color{muted} #3}
  \end{center}
  \vspace{0.6em}}

\fancyhead[L]{\small\color{muted}集成投票的极限：初学者导读}
\fancyhead[R]{\small\color{muted}Part I \& Part II}
\hypersetup{pdftitle={集成投票的极限：Part I 与 Part II 证明导读}}
\newcommand{\ok}{$\checkmark$}
\newcommand{\no}{\textcolor{warm}{$\times$}}

\begin{document}
\doctitle{集成投票的极限}{Part I 与 Part II 证明导读（面向初学者）}{2026 年 10 月}

\begin{keybox}[title=这份讲义讲什么]
很多人相信：“把几个还不错的模型放在一起投票，结果一定更好。”这份讲义用三个问题检验这个说法：
\begin{enumerate}[leftmargin=1.6em]
\item 如果各个模型\textbf{互相独立}，投票确实越来越好吗？（会，这是 Condorcet 陪审团定理。）
\item 如果模型之间可以\textbf{任意相关}，只知道每个模型的准确率，投票最差能差到什么程度？
\item 如果再告诉你\textbf{每两个模型之间的相关系数}，能不能把投票准确率算准？（不能，差距可能很大。）
\end{enumerate}
\textbf{需要的预备知识：}概率、期望、二项分布、相关系数的定义。每节末尾有“要点”框，可以先看要点再读细节。严格的定理与证明见配套文档《严格数学证明》。
\end{keybox}

\tableofcontents

%=====================================================================
\section{把问题翻译成数学}
%=====================================================================

设想三个分类器判断邮件是否为垃圾邮件。对每封测试邮件，我们只记录每个分类器\textbf{判对}（记作 1）还是\textbf{判错}（记作 0）：

\begin{center}
\begin{tabular}{cccc|c|c}
\toprule
邮件 & 分类器 1 & 分类器 2 & 分类器 3 & 判对的个数 $S$ & 多数投票对吗？\\
\midrule
1 & 1 & 1 & 0 & 2 & 对\\
2 & 0 & 1 & 0 & 1 & 错\\
3 & 1 & 1 & 1 & 3 & 对\\
4 & 0 & 0 & 1 & 1 & 错\\
\bottomrule
\end{tabular}
\end{center}

\paragraph{为什么只看“对或错”就够了？}这是二分类问题，只有两个标签。判对的分类器都投给真实标签，判错的都投给另一个标签。所以真实标签得到 $S$ 票，错误标签得到 $3-S$ 票。多数投票判对，当且仅当 $S\ge2$。一般地，$K$ 个分类器（$K$ 为奇数）投票时，多数投票判对当且仅当至少有
\[
m=\frac{K+1}{2}
\]
个分类器判对。

于是我们定义：$Z_i=1$ 表示第 $i$ 个分类器判对，$p_i=\Prob(Z_i=1)$ 是它的准确率，$S=Z_1+\dots+Z_K$ 是判对的个数。整个问题的主角是
\[
\boxed{\ \text{集成准确率}\ \ A=\Prob(S\ge m).\ }
\]

\begin{keybox}[title=要点]
集成准确率只取决于“哪些分类器同时判对”的规律，也就是 $(Z_1,\dots,Z_K)$ 的联合分布。不同的联合分布可以有完全相同的单个准确率，但 $A$ 不同。这是整份讲义的核心。
\end{keybox}

%=====================================================================
\section{独立时：分类器越多越好}
%=====================================================================

\subsection{三个分类器的计算}
如果三个分类器互相独立、准确率都是 $p=0.7$，那么 $S\sim\Bin(3,0.7)$：
\[
A=\Prob(S=2)+\Prob(S=3)=3\times0.7^2\times0.3+0.7^3=0.441+0.343=0.784 .
\]
三个 70\% 的分类器投票，得到了 78.4\% 的准确率。

\subsection{为什么多加两个会更好？}
从 $K$ 个分类器加到 $K+2$ 个，多数门槛从 $m$ 变成 $m+1$。新来的两票只在两种“临界”局面下才会改变结果：

\begin{center}
\begin{tabular}{lll}
\toprule
原来 $K$ 个中判对的个数 & 原来的结果 & 加两个之后\\
\midrule
差一票（$S=m-1$） & 错 & 两个新的都对 $\Rightarrow$ 变成对（概率 $p^2$）\\
刚好过半（$S=m$） & 对 & 两个新的都错 $\Rightarrow$ 变成错（概率 $(1-p)^2$）\\
其他情况 & — & 不变\\
\bottomrule
\end{tabular}
\end{center}

独立时，可以算出“差一票”的概率恰好是“刚好过半”概率的 $\frac{1-p}{p}$ 倍。所以
\[
\text{增量}=\Prob(\text{刚好过半})\Bigl[\tfrac{1-p}{p}\,p^2-(1-p)^2\Bigr]=\Prob(\text{刚好过半})\,(1-p)\,(2p-1).
\]
只要 $p>1/2$，增量就是正的。例如 $p=0.7$ 时从 3 个加到 5 个，增量为 $3\times0.7^2\times0.3^2\times0.4\approx0.053$，准确率从 0.784 升到 0.837。

\subsection{分类器无穷多时}
由大数定律，$S/K$ 会集中在 $p$ 附近。$p>1/2$ 时，“判对的超过一半”以趋于 1 的概率发生，于是 $A\to1$。Hoeffding 不等式还给出了速度：$1-A\le e^{-2K(p-1/2)^2}$。$p<1/2$ 时则相反，$A\to0$。

\begin{keybox}[title=要点]
独立且准确率都相同时，只要每个分类器都好于随机猜测（$p>1/2$），加分类器就会让投票准确率单调上升，并趋于 1。关键在于\textbf{独立}这个假设。（准确率不同时未必如此：一个 0.99 的分类器加上两个独立的 0.6 分类器，投票准确率只有 0.835。）
\end{keybox}

%=====================================================================
\section{不再假设独立：最坏能有多坏？}
%=====================================================================

\subsection{十封邮件的例子}
三个分类器，每个都恰好判对 10 封中的 6 封（$p=0.6$）。下面两种安排都满足这一点，但投票结果相差很大。

\begin{center}
\small
\begin{tabular}{c|ccc|c}
\multicolumn{5}{c}{\textbf{安排 A（最坏）}}\\
\toprule
邮件 & C1 & C2 & C3 & 投票\\
\midrule
1--4 & \ok & \ok & \ok & 对\\
5, 6 & \ok & \no & \no & 错\\
7, 8 & \no & \ok & \no & 错\\
9, 10 & \no & \no & \ok & 错\\
\midrule
\multicolumn{4}{l|}{投票准确率} & \textbf{0.4}\\
\bottomrule
\end{tabular}
\qquad
\begin{tabular}{c|ccc|c}
\multicolumn{5}{c}{\textbf{安排 B（最好）}}\\
\toprule
邮件 & C1 & C2 & C3 & 投票\\
\midrule
1--3 & \ok & \ok & \no & 对\\
4--6 & \ok & \no & \ok & 对\\
7--9 & \no & \ok & \ok & 对\\
10 & \no & \no & \no & 错\\
\midrule
\multicolumn{4}{l|}{投票准确率} & \textbf{0.9}\\
\bottomrule
\end{tabular}
\end{center}

在安排 A 中，每个分类器都判对 $4+2=6$ 封；在安排 B 中，每个分类器都判对 $3+3=6$ 封。单看每个分类器，两种安排完全一样，但投票准确率一个是 40\%（比单个分类器还差），一个是 90\%。独立时是 64.8\%，介于两者之间。

\subsection{“数票”论证：为什么最坏是 0.4}
关键想法是数“判对的票”的总数。

\begin{itemize}
\item 10 封邮件 $\times$ 3 个分类器 $\times$ 0.6 = \textbf{18 张判对票}，这个总数与分类器之间是否相关无关。
\item 投票\textbf{失败}的邮件最多只有 1 张判对票（否则就过半了）。
\item 投票\textbf{成功}的邮件最多有 3 张判对票。
\end{itemize}
设有 $x$ 封成功，则 $18\le 3x+1\cdot(10-x)$，解得 $x\ge4$。所以无论怎样安排，准确率至少是 40\%。安排 A 正好把这个不等式变成了等式：成功的邮件“满票”通过，失败的邮件每封都只差一票。

同理，成功的邮件\emph{至少}要 2 张判对票，所以 $18\ge2x$，即 $x\le9$。最好是 90\%，安排 B 正好取到：每封成功邮件都“刚好过半”，一票都不浪费。

\begin{goodbox}[title=一般公式（$K$ 个分类器，准确率都是 $p$）]
把上面的数票推广：平均每个样本有 $Kp$ 张判对票；失败的样本最多 $m-1$ 张，成功的样本最多 $K$ 张，至少 $m$ 张。于是
\[
\Amin=\max\Bigl\{0,\ \frac{2Kp-K+1}{K+1}\Bigr\},\qquad \Aplus=\min\Bigl\{1,\ \frac{2Kp}{K+1}\Bigr\}.
\]
这两个界都能取到。最坏情况（当它大于 0 时）用“要么全对，要么恰好差一票”的安排；最好情况（当它小于 1 时）用“要么刚好过半，要么全错”的安排。
\end{goodbox}

\subsection{分类器越多，最坏情况越糟}
把 $p=0.7$ 代入，比较独立情形和最坏情形：

\begin{center}
\begin{tikzpicture}
\begin{axis}[width=0.82\linewidth,height=5.6cm,xlabel={分类器个数 $K$},ylabel={投票准确率},
  xmin=0,xmax=22,ymin=0.38,ymax=1.03,xtick={1,3,5,7,9,11,13,15,17,19,21},grid=major,grid style={gray!20},
  legend style={at={(0.98,0.5)},anchor=east,font=\small,draw=gray!40,fill=white},tick label style={font=\small},label style={font=\small}]
\addplot[accent,very thick,mark=*,mark size=1.6pt] coordinates {(1,0.7000)(3,0.7840)(5,0.8369)(7,0.8740)(9,0.9012)(11,0.9218)(13,0.9376)(15,0.9500)(17,0.9597)(19,0.9674)(21,0.9736)};
\addlegendentry{独立}
\addplot[warm,very thick,mark=square*,mark size=1.6pt] coordinates {(1,0.7000)(3,0.5500)(5,0.5000)(7,0.4750)(9,0.4600)(11,0.4500)(13,0.4429)(15,0.4375)(17,0.4333)(19,0.4300)(21,0.4273)};
\addlegendentry{最坏情况}
\addplot[gray,thin,dotted] coordinates {(0,0.4)(22,0.4)};
\node[font=\scriptsize,gray,anchor=south west] at (axis cs:1.3,0.40) {极限 $2p-1=0.4$};
\end{axis}
\end{tikzpicture}
\end{center}

当 $p>1/2$ 时，最坏情况可以写成 $\Amin=(2p-1)+\frac{2(1-p)}{K+1}$，随 $K$ 增大而\emph{下降}，极限是 $2p-1$。$p=0.7$ 时，极限只有 0.4，远低于单个分类器的 0.7。

\begin{warnbox}[title=这对实践意味着什么]
最坏情况并不是“所有分类器犯同样的错”：完全相同的分类器投票，准确率就是 $p$，既没变好也没变差。最坏情况是一种更“刁钻”的相关结构：在一部分样本上全部判对（多出来的正确票被浪费），在其余样本上恰好差一票而失败。这说明评估集成时不能只看单个准确率，还要看错误在样本上的分布方式。集成的收益来自错误的\textbf{多样性}，而不仅仅是模型的数量。
\end{warnbox}

\begin{keybox}[title=要点]
只知道单个准确率时，集成准确率可以落在一个区间里的任何位置。区间的两端由“数票”不等式给出，而且都能取到。分类器越多，最坏情况反而越差。
\end{keybox}

%=====================================================================
\section{准确率各不相同时}
%=====================================================================

\subsection{一个启发性的例子}
设三个分类器的准确率为 $(0.95,\,0.95,\,0.3)$。如果直接套用“所有分类器一起数票”，平均每个样本有 $2.2$ 张判对票，得到 $A\ge(2.2-1)/2=0.6$。但其实可以得到更好的下界：只看前两个强分类器，它们同时判对的概率至少是 $0.95+0.95-1=0.9$；而只要它们同时判对，多数投票就一定对。所以 $A\ge0.9$，而且这正是最坏情况。

启示：数票时不一定要把所有分类器都算进去，挑一组最强的分类器来数可能更好。

\subsection{计算步骤}
\begin{goodbox}[title=最坏情况的计算方法]
\begin{enumerate}[leftmargin=1.6em]
\item 把准确率从大到小排好：$\pd1\ge\pd2\ge\dots\ge\pd K$。
\item 对每个 $j=m,m+1,\dots,K$，只对最强的 $j$ 个分类器数票，得到
\[
c_j=\frac{\pd1+\dots+\pd j-(m-1)}{j-m+1}.
\]
\item 最坏准确率 $\Amin=\max\{0,c_m,c_{m+1},\dots,c_K\}$。
\end{enumerate}
\end{goodbox}

\textbf{例：}$K=5$，$m=3$，准确率 $(0.9,0.8,0.7,0.6,0.5)$。
\[
c_3=\frac{2.4-2}{1}=0.4,\qquad c_4=\frac{3.0-2}{2}=0.5,\qquad c_5=\frac{3.5-2}{3}=0.5 .
\]
所以 $\Amin=0.5$。最好情况有对称的公式：把准确率从小到大排序，
\[
\Aplus=\min\Bigl\{1,\ \min_{m\le j\le K}\frac{\pa1+\dots+\pa j}{j-m+1}\Bigr\}.
\]
这个例子的最好情况是 1。

\subsection{为什么这个界能取到：梳子构造}
要证明“最坏真的能这么坏”，需要构造一个具体的联合分布。核心工具是一个巧妙的随机构造，叫\textbf{系统抽样}，可以想象成一把梳子：

\begin{center}
\begin{tikzpicture}[x=3.2cm,y=1cm]
\def\U{0.35}
\fill[accent!25] (0,0) rectangle (0.9,0.5);
\fill[good!25] (0.9,0) rectangle (1.7,0.5);
\fill[warm!25] (1.7,0) rectangle (2.4,0.5);
\fill[accent!12] (2.4,0) rectangle (3.0,0.5);
\fill[gray!25] (3.0,0) rectangle (3.5,0.5);
\foreach \a/\b/\n/\v in {0/0.9/1/0.9,0.9/1.7/2/0.8,1.7/2.4/3/0.7,2.4/3.0/4/0.6,3.0/3.5/5/0.5}{
  \draw (\a,0) rectangle (\b,0.5);
  \node[font=\small] at ({(\a+\b)/2},0.25) {C\n};
  \node[font=\scriptsize,muted] at ({(\a+\b)/2},-0.25) {\v};
}
\draw[-{Stealth}] (-0.1,0) -- (3.75,0);
\foreach \x in {0,1,2,3} \node[font=\scriptsize,below] at (\x,-0.45) {\x};
\foreach \k in {0,1,2,3}{
  \draw[very thick,red!70!black] ({\U+\k},1.25) -- ({\U+\k},0.55);
  \fill[red!70!black] ({\U+\k},0.55) circle (1.5pt);
}
\draw[very thick,red!70!black] (\U,1.25) -- ({\U+3},1.25);
\node[font=\small,red!70!black,anchor=west] at ({\U+3.05},1.25) {梳子（齿距 1，随机平移 $U$）};
\node[font=\scriptsize,anchor=south] at (\U,1.3) {$U$};
\end{tikzpicture}
\end{center}

\begin{enumerate}
\item 把长度为 $p_1,p_2,\dots,p_K$ 的小段首尾相接，排在数轴上。
\item 取一把齿距为 1 的梳子，随机平移一个 $U\sim\mathrm{Unif}[0,1)$ 后放上去。
\item 哪一段被梳齿扎到，对应的分类器就“判对”。
\end{enumerate}

这个构造有两个好性质：
\begin{itemize}
\item 每段长度不超过 1，所以最多被扎一次；被扎到的概率恰好等于段长，即每个分类器的准确率恰好是 $p_i$。
\item 被扎到的段数，也就是判对的个数，永远是 $\lfloor\sum p_i\rfloor$ 或 $\lceil\sum p_i\rceil$，几乎没有波动。
\end{itemize}
上图中总长 $3.5$，所以每次都有 3 个或 4 个分类器判对；图中这一次梳齿扎到 C1、C2、C3、C5，即 4 个分类器判对。最坏情况的构造是把两把梳子按比例混合，但两把梳子用的段长要先改一改：一把的段长是 $u_i=(p_i-a)_+/(1-a)$，总长不超过 $m-1$，保证判对数 $\le m-1$（失败）；另一把的段长是 $v_i=\min\{p_i,a\}/a$，总长至少 $m$，保证判对数 $\ge m$（成功）。这里 $a=\Amin$ 是混合比例，混合后每个分类器的准确率恰好回到 $p_i$。以上面的 $K=5$ 为例，$a=0.5$：以 $\frac12$ 的概率五个全对（$v$ 全为 1）；以 $\frac12$ 的概率用段长 $(0.8,0.6,0.4,0.2,0)$ 的梳子，总长恰为 2，所以恰好 2 个判对，投票失败。可以验证分类器 1 的准确率为 $\frac12\times1+\frac12\times0.8=0.9$。

\begin{keybox}[title=要点]
准确率不同时，最坏情况等于“挑一组最强的 $j$ 个分类器数票”得到的最好下界。这个下界同样能取到，构造工具是梳子式的系统抽样。
\end{keybox}

%=====================================================================
\section{Part II：知道两两相关系数也不够}
%=====================================================================

\subsection{相关系数对 0/1 变量意味着什么}
对两个 0/1 变量，$\E[Z_iZ_j]=\Prob(\text{两个都判对})$。由相关系数的定义，
\[
r:=\Prob(\text{分类器 } i,j \text{ 都判对})=p^2+\rho\,p(1-p).
\]
所以知道 $p$ 和 $\rho$，就等于知道了“任意两个同时判对”的概率 $r$。例如 $\rho=0$ 时 $r=p^2$。

\subsection{关键恒等式}
三个分类器时，多数投票判对的指示函数可以写成
\[
\1\{S\ge2\}=Z_1Z_2+Z_1Z_3+Z_2Z_3-2Z_1Z_2Z_3 .
\]
逐一验证：

\begin{center}
\begin{tabular}{cccc}
\toprule
判对个数 $S$ & 两两乘积之和 & $-2Z_1Z_2Z_3$ & 合计\\
\midrule
0 或 1 & 0 & 0 & 0\\
2 & 1 & 0 & 1\\
3 & 3 & $-2$ & 1\\
\bottomrule
\end{tabular}
\end{center}

取期望，得到
\[
\boxed{\ A=3r-2t,\qquad t=\Prob(\text{三个都判对}).\ }
\]
$p$ 和 $\rho$ 决定了 $r$，却\textbf{没有}决定 $t$。$t$ 是一个“三阶”信息，只看两两关系看不到它。这就是相关系数不够用的根本原因。

\subsection{25 封邮件的例子}
取 $p=0.6$，$\rho=0$（两两不相关），于是 $r=0.36$。下面是两种安排，以及独立情形（独立情形需要 125 封邮件才能凑成整数）：

\begin{center}
\small
\begin{tabular}{lccc}
\toprule
判对模式 & 最坏（25 封） & 独立（125 封） & 最好（25 封）\\
\midrule
三个都对 & 7 & 27 & 3\\
恰好两个对（三种模式，每种） & 2 & 18 & 6\\
恰好一个对（三种模式，每种） & 4 & 12 & 0\\
都不对 & 0 & 8 & 4\\
\midrule
每个分类器判对的比例 & $15/25=0.6$ & $75/125=0.6$ & $15/25=0.6$\\
每两个同时判对的比例 & $9/25=0.36$ & $45/125=0.36$ & $9/25=0.36$\\
\midrule
\textbf{投票准确率} & $\mathbf{13/25=0.52}$ & $\mathbf{81/125=0.648}$ & $\mathbf{21/25=0.84}$\\
\bottomrule
\end{tabular}
\end{center}

以最坏情况为例核对：分类器 1 判对的邮件有“三个都对”7 封、含它的两种“两个对”模式各 2 封、“只有它对”4 封，共 $7+2+2+4=15$ 封。分类器 1 和 2 同时判对的有 $7+2=9$ 封，比例 $0.36=0.6^2$，所以相关系数为 0。

三种安排的单个准确率和两两相关系数完全相同，投票准确率却从 52\% 到 84\% 不等。

\begin{warnbox}[title=常见误解]
“两两不相关”不等于“相互独立”。上表最坏和最好的两列都是两两不相关的，但都不独立。只有中间那一列才是真正的独立。
\end{warnbox}

\subsection{$t$ 能取哪些值？}
用 $q_s$ 表示恰好 $s$ 个判对的比例。由“总比例为 1”“平均判对数为 $3p$”“平均判对对数为 $3r$”三个条件，可以把所有 $q_s$ 用 $t$ 表示：
\[
q_0=1-3p+3r-t,\quad q_1=3p-6r+3t,\quad q_2=3r-3t,\quad q_3=t .
\]
比例不能是负数，这给出 $t$ 的范围 $\max\{0,2r-p\}\le t\le\min\{r,1-3p+3r\}$。代入 $A=3r-2t$：

\begin{goodbox}[title=三个分类器的精确区间]
\[
A_{\min}=\max\{r,\ 6p-3r-2\},\qquad A_{\max}=\min\{3r,\ 2p-r\},\qquad r=p^2+\rho p(1-p).
\]
\end{goodbox}

直观上，最坏情况让“三个都对”尽量多（把判对票浪费在已经稳赢的样本上），最好情况让“三个都对”尽量少。

\paragraph{不是所有的 $\rho$ 都可能。}比如三枚“公平硬币”（$p=0.5$）不可能两两都强烈负相关：如果 1 和 2 总是相反，1 和 3 也总是相反，那么 2 和 3 就总是相同。由 $q_s\ge0$ 可以算出，$p=0.5$ 时 $\rho$ 至少是 $-\frac13$。这比相关矩阵半正定给出的 $-\frac12$ 更严格。

\paragraph{准确率各不相同时}做法完全一样：8 种判对模式的概率都可以用 $t$ 表示，令它们非负就得到 $t$ 的范围 $[L,U]$，于是 $A\in[R-2U,R-2L]$，其中 $R$ 是三个“两两同时判对”概率之和。

\begin{keybox}[title=要点]
准确率加上全部两两相关系数，仍然不能唯一确定投票准确率，因为结果还依赖“三个同时判对”的概率。三个分类器时，这个不确定性可以精确地写成一个区间。
\end{keybox}

%=====================================================================
\section{更多分类器：交给线性规划}
%=====================================================================

$K$ 个分类器有 $2^K$ 种判对模式。把每种模式的概率当作未知数，那么
\begin{itemize}
\item 投票准确率是“过半模式”的概率之和，是未知数的\textbf{线性}函数；
\item “概率之和为 1”“每个分类器的准确率”“每两个同时判对的概率”都是\textbf{线性}等式；
\item 概率非负是线性不等式。
\end{itemize}
所以，求投票准确率的最小值和最大值就是一个\hspace{0pt}\textbf{线性规划}\hspace{0pt}问题，可以用标准工具数值求解，例如 Python 的
\texttt{scipy.optimize.linprog}。

如果所有分类器的准确率相同、两两相关系数也相同，问题还能化简：只需要知道 $S$ 的分布 $q_0,\dots,q_K$，未知数从 $2^K$ 个降到 $K+1$ 个。下表是 $p=0.6$、$\rho=0$ 时的结果：

\begin{center}
\small
\begin{tabular}{cccc}
\toprule
$K$ & 最坏 & 独立 & 最好\\
\midrule
3 & 0.520 & 0.648 & 0.840\\
5 & \textbf{0.467} & 0.683 & 0.920\\
7 & 0.480 & 0.710 & 0.960\\
11 & 0.500 & 0.753 & 1.000\\
21 & 0.580 & 0.826 & 1.000\\
\bottomrule
\end{tabular}
\end{center}

注意 $K=5$ 那一行：五个两两\emph{不相关}、准确率 60\% 的分类器，多数投票的准确率最低可以只有 46.7\%，比单个分类器还差。

%=====================================================================
\section{总结}
%=====================================================================

\begin{center}
\small
\renewcommand{\arraystretch}{1.3}
\begin{tabular}{>{\raggedright\arraybackslash}p{0.27\linewidth}>{\raggedright\arraybackslash}p{0.62\linewidth}}
\toprule
已知的信息 & 能说什么\\
\midrule
各自准确率 + 独立 & 准确率可以精确算出；准确率都为 $p>1/2$ 时越多越好，趋于 1\\
只有各自准确率 & 只能确定一个区间；最坏情况可以比单个分类器还差，且随 $K$ 变差\\
准确率 + 两两相关 & 区间变窄，但通常仍不是一个点；$K=3$ 时有闭式区间\\
完整的联合分布 & 准确率被唯一确定\\
\bottomrule
\end{tabular}
\end{center}

\begin{warnbox}[title=四个常见误区]
\begin{enumerate}[leftmargin=1.6em]
\item \textbf{“不相关就是独立。”}对 0/1 变量，两两不相关不意味着相互独立（见第 5 节的表格）。
\item \textbf{“多加模型一定更好。”}在准确率相同且独立时成立；依赖结构不利时，最坏情况下 $K$ 越大越差。
\item \textbf{“知道相关系数就能算出集成准确率。”}还需要三阶及更高阶的信息。
\item \textbf{“最坏情况只是一个保守的估计。”}这里的最坏情况是紧的，能被具体的联合分布达到。
\end{enumerate}
\end{warnbox}

\paragraph{这些结论从哪里来？}独立情形是 18 世纪的 Condorcet 陪审团定理。只知道准确率时的最坏与最好情况，在分类器集成文献中由 Kuncheva 等人（2003）和 Matan（1996）给出，在陪审团理论中由 Zaigraev 与 Kaniovski（2012）给出。给定两两相关系数时用线性规划求上下界的思路见 Kaniovski（2010）。配套文档《引用已有结论的简明证明》给出了完整的引用。

\end{document}
